\begingroup
\renewcommand{\endProof}{\par}
\begin{Proof}
设 $h\in(0,R(1-\lambda)/2)$, 固定指标 $k=1,\ldots,n-1$ (即 $k\ne n$). 选取 $\zeta\in C_b^\infty(D^+)$, 使得\begin{enumerate}\item $0\leq\zeta\leq1$,\item 在 $Q^+$ 上 $\zeta\equiv1$,\item $U:=\operatorname{supp}\zeta\subset B_{R(1+\lambda)/2}(0)$.\end{enumerate}\textbf{Note:} 函数 $\zeta$ 不属于 $\mathcal{D}(D^+)$, 因为在 $D^+$ 边界的平坦部分 $\{x_n=0\}$ 上 $\zeta\not\equiv0$.

现在定义
\begin{equation}\label{eq:flat-boundary-tangential-difference-test-function}
v:=-D_k^{-h}(\zeta^2D_k^hu).
\end{equation}
利用差商的定义作一些运算后, 得到如下恒等式.
\begin{equation}\label{eq:flat-boundary-test-function-translation-expansion}
\begin{aligned}
v(\bm{x})=-\frac{1}{h^2}\Bigl(&\zeta^2(\bm{x})[u(\bm{x}+h\bm{e}_k)-u(\bm{x})]\\
&+\zeta^2(\bm{x}-h\bm{e}_k)[u(\bm{x}-h\bm{e}_k)-u(\bm{x})]\Bigr).
\end{aligned}
\end{equation}
注意构造 $v$ 时, 只沿平面 $x_n=0$ 的切向作了平移. 关键想法是可以沿平面 $x_n=0$ ``滑动 $u$ 的支集\textquotesingle\textquotesingle, 而不破坏边界条件. 还应注意, 这些平移均不会将 $\zeta$ 的支集移出 $\overline{D^+}$. 这些事实保证\begin{enumerate}\item $v$ 在 $D^+$ 上有良好定义,\item $v\in H^1(D^+)$ (因为 $u\in H^1(D^+)$),\item $v\in H_0^1(D^+)$ (因为 $\zeta$ 在 $D^+$ 边界的弯曲部分为零, $u$ 在边界的平坦部分为零, 而 $\zeta$ 和 $u$ 的所有这些平移也一样).\end{enumerate}定义
\begin{equation}\label{eq:flat-boundary-localized-tangential-difference}
w:=\zeta D_k^hu,
\end{equation}
从而 $v=-D_k^{-h}(\zeta w)$.

现在按照与 Theorem~\ref{thm:second-order-interior-regularity} 中估计 \eqref{eq:interior-regularity-shifted-principal-form-bound} 的推导类似的过程进行, 但使用上面定义的 $v$. 得到\[\begin{aligned}
-\int_{D^+}gD_k^{-h}(\zeta w)\,\mathrm{d}\bm{x}
&=\int_{D^+}a_{ij}u_{x_j}v_{x_i}\,\mathrm{d}\bm{x}\\
&=\int_{D^+}D_k^h(a_{ij}u_{x_j})(\zeta w)_{x_i}\,\mathrm{d}\bm{x}\\
&=\int_{D^+}a_{ij}(\bm{x}+h\bm{e}_k)D_k^hu_{x_j}(\zeta w)_{x_i}\,\mathrm{d}\bm{x}\\
&\quad+\int_{D^+}D_k^ha_{ij}u_{x_j}(\zeta w)_{x_i}\,\mathrm{d}\bm{x}\\
&=\int_{D^+}a_{ij}(\bm{x}+h\bm{e}_k)w_{x_j}w_{x_i}\,\mathrm{d}\bm{x}\\
&\quad-\int_{D^+}a_{ij}(\bm{x}+h\bm{e}_k)\zeta_{x_j}D_k^huw_{x_i}\,\mathrm{d}\bm{x}\\
&\quad+\int_{D^+}a_{ij}(\bm{x}+h\bm{e}_k)D_k^hu_{x_j}\zeta_{x_i}w\,\mathrm{d}\bm{x}\\
&\quad+\int_{D^+}D_k^ha_{ij}u_{x_j}(\zeta w)_{x_i}\,\mathrm{d}\bm{x}.
\end{aligned}\]这里 $g$ 如 \eqref{eq:interior-regularity-reduced-right-hand-side} 定义. 整理得到
\begin{equation}\label{eq:flat-boundary-principal-energy-four-term-decomposition}
\int_{D^+}a_{ij}(\bm{x}+h\bm{e}_k)w_{x_j}w_{x_i}\,\mathrm{d}\bm{x}=I_1+I_2+I_3+I_4,
\end{equation}
其中
\begin{equation}\label{eq:flat-boundary-energy-term-one}
I_1:=-\int_{D^+}gD_k^{-h}(\zeta w)\,\mathrm{d}\bm{x},
\end{equation}

\begin{equation}\label{eq:flat-boundary-energy-term-two}
I_2:=\int_{D^+}a_{ij}(\bm{x}+h\bm{e}_k)\zeta_{x_j}D_k^huw_{x_i}\,\mathrm{d}\bm{x},
\end{equation}

\begin{equation}\label{eq:flat-boundary-energy-term-three}
I_3:=-\int_{D^+}a_{ij}(\bm{x}+h\bm{e}_k)D_k^hu_{x_j}\zeta_{x_i}w\,\mathrm{d}\bm{x},
\end{equation}

\begin{equation}\label{eq:flat-boundary-energy-term-four}
I_4:=-\int_{D^+}D_k^ha_{ij}u_{x_j}(\zeta w)_{x_i}\,\mathrm{d}\bm{x}.
\end{equation}
利用读者从 Theorem~\ref{thm:second-order-interior-regularity} 的证明中应已熟悉的技巧 (主要是 Hölder 不等式和 Lemma~\ref{lem:difference-quotient-weak-derivative-bound}) 估计这些项. 将利用估计
\begin{equation}\label{eq:flat-boundary-test-function-l2-bound}
\|v\|_2=\|D_k^{-h}(\zeta w)\|_2\leq2\|\zeta\|_{1,\infty}\|w\|_{1,2}.
\end{equation}

\end{Proof}
\endgroup
